\section{Техническое задание}

\subsection{Основание для разработки}

Полное наименование системы: «Разработка компонентно-ориентированной программной системы для построения интерактивных приложений на Java».

Основанием для разработки программы является задание на выпускную квалификационную работу бакалавра по направлению подготовки 09.03.04 «Программная инженерия». Разрабатываемая программная система представляет собой программную основу для создания интерактивных приложений на языке Java. В дальнейшем по мере расширения подсистем, средств разработки и пользовательской документации данная система может быть развита до более полного игрового движка, однако в рамках настоящей работы основной акцент сделан на ядре, компонентной модели, сценах, системах обработки, визуализации, вводе и базовой физике.

Разработка ведется с учетом требований к стадиям создания программной документации. В разделе технического задания необходимо определить назначение разработки, сформулировать задачи, описать требования к программной системе и показать порядок ее использования разработчиком интерактивного приложения.

\subsection{Цель и назначение разработки}

Целью разработки является создание компонентно-ориентированной программной системы, позволяющей Java-разработчику строить простые интерактивные приложения и учебные игровые прототипы без самостоятельной реализации полного набора низкоуровневых механизмов: создания окна, основного цикла приложения, управления сценой, хранения объектов, обработки пользовательского ввода, визуализации трехмерных объектов и базовой обработки физических тел.

Функциональное назначение системы состоит в предоставлении разработчику набора классов и подсистем, с помощью которых можно описывать интерактивную сцену как совокупность сущностей, компонентов и систем обработки. Сущность используется как контейнер компонентов, компонент хранит состояние или поведение объекта, а система выполняет обработку групп сущностей, удовлетворяющих заданным условиям. Такой подход позволяет отделить общие механизмы работы приложения от прикладной логики конкретной сцены.

Эксплуатационное назначение системы заключается в использовании ее Java-разработчиком при создании настольных интерактивных приложений. Пользователь разрабатываемой системы не является конечным пользователем готовой игры или демонстрации. В данном случае пользователем выступает программист, который подключает исходные коды движка, создает класс приложения, описывает сцену, добавляет объекты, настраивает компоненты и запускает приложение в среде разработки или из командной строки.

Для достижения поставленной цели необходимо решить следующие задачи:
{
\renewcommand{\labelenumi}{\arabic{enumi}.}
\begin{enumerate}
\item Разработать ядро компонентно-ориентированной модели, включающее сущности, компоненты, системы обработки, запросы сущностей и жизненный цикл объектов сцены.

\item Спроектировать и реализовать подсистему управления приложением и сценами, обеспечивающую создание окна, выполнение основного цикла, обработку времени кадра и безопасное изменение состава сцены во время работы программы.

\item Реализовать графическую подсистему на основе OpenGL, обеспечивающую отображение объектов сцены с использованием мешей, материалов, текстур, камер и направленного освещения.

\item Реализовать подсистемы пользовательского ввода, скриптов и базовой физики, позволяющие обрабатывать действия пользователя, задавать поведение объектов средствами Java-кода и демонстрировать взаимодействие физических тел и коллайдеров.

\item Подготовить набор демонстрационных сцен и описание порядка использования программной системы с примерами кода, показывающими создание приложения, сцены, сущностей, компонентов, систем обработки и пользовательских скриптов.
\end{enumerate}
}

\subsection{Требования к программной системе}

Разрабатываемая программная система должна предоставлять Java-разработчику средства для создания интерактивного приложения программным способом. В системе не предусматривается визуальный редактор сцен, поэтому создание объектов, настройка компонентов, подключение систем и запуск сцены выполняются в исходном коде приложения.

Основным сценарием использования системы является разработка интерактивной сцены. Разработчик создает объект приложения, задает параметры окружения, формирует сцену, добавляет в нее системы обработки, затем создает сущности с компонентами и запускает основной цикл выполнения. После запуска система самостоятельно выполняет обновление активной сцены, обработку ввода, вызов пользовательских скриптов, физическое обновление, рендеринг и освобождение ресурсов при завершении работы.

\subsubsection{Требования к архитектурной модели}

Архитектурная модель системы должна основываться на компонентно-ориентированном подходе. В системе должны использоваться следующие основные понятия:
\begin{itemize}
\item сущность -- объект сцены, имеющий уникальный идентификатор и набор компонентов;
\item компонент -- часть состояния или поведения сущности;
\item система -- объект, выполняющий обработку сущностей, имеющих требуемый набор компонентов;
\item сцена -- контейнер сущностей и систем, обновляемый в основном цикле приложения;
\item контекст движка -- объект, объединяющий окно, ввод, менеджер ресурсов, менеджер сцен, конфигурацию и время выполнения.
\end{itemize}

Компонентная модель должна позволять добавлять к сущности несколько независимых возможностей. Например, объект сцены может одновременно иметь компонент трансформации, компонент отображения, физическое тело и коллайдер. При этом прикладная логика не должна требовать создания отдельного класса-наследника для каждого нового типа игрового объекта.

Система должна поддерживать жизненный цикл компонентов. После присоединения компонента к сущности должен вызываться этап инициализации связи с сущностью, при запуске сцены -- этап начала работы, при остановке -- этап остановки, при удалении -- этап уничтожения. Это необходимо для корректной регистрации внутренних состояний, освобождения ресурсов и прекращения работы скриптов.

\subsubsection{Требования к управлению приложением и сценами}

Программная система должна обеспечивать создание окна приложения, настройку базовых параметров запуска и выполнение основного цикла. К базовым параметрам относятся цвет очистки окна, масштаб времени, фиксированный шаг обновления, автоматическое обновление соотношения сторон камеры и режим отладочной проверки ошибок OpenGL.

Менеджер сцен должен обеспечивать загрузку активной сцены и замену текущей сцены во время выполнения. Смена сцены должна выполняться безопасно, без нарушения текущего цикла обновления. Это важно для интерактивных приложений, в которых сцена может быть перезапущена по действию пользователя, например при нажатии клавиши сброса демонстрации.

Сцена должна поддерживать добавление и удаление сущностей, добавление и удаление систем обработки, выполнение цикла обновления и освобождение ресурсов. Изменение структуры сцены во время обновления должно откладываться до безопасного момента, чтобы добавление или удаление объектов из пользовательского скрипта не приводило к ошибкам обхода коллекций.

\subsubsection{Требования к графической подсистеме}

Графическая подсистема должна использовать OpenGL и обеспечивать отображение трехмерных объектов в окне приложения. Для типового объекта сцены необходимо поддерживать компонент трансформации, компонент отображения меша и материал. Меш должен задавать геометрию объекта, материал -- параметры отображения, а трансформация -- положение, поворот и масштаб в сцене.

В составе системы должны быть предусмотрены стандартные геометрические примитивы, необходимые для демонстрационных сцен. К таким примитивам относятся куб, плоскость или четырехугольник, а также простые формы, позволяющие быстро собрать тестовую сцену без загрузки модели из внешнего файла. Загрузка мешей из внешних файлов в рамках текущей версии системы не является обязательной возможностью.

Подсистема ресурсов должна обеспечивать загрузку и повторное использование шейдеров и текстур. Текстуры должны загружаться из файлов ресурсов с использованием средств LWJGL STB. Повторная загрузка уже использованного ресурса должна избегаться за счет кэширования, так как в интерактивных приложениях один и тот же шейдер или одна и та же текстура могут использоваться многими объектами.

Для освещения должна быть предусмотрена поддержка базового направленного источника света. Такая модель освещения достаточна для демонстрации объемных объектов, положения камеры и работы материалов, но не рассматривается как полноценная система освещения с разными типами источников, тенями и постобработкой.

\subsubsection{Требования к пользовательскому вводу}

Программная система должна обеспечивать получение состояния клавиатуры и мыши. Разработчик должен иметь возможность проверять факт удержания клавиши, однократное нажатие клавиши, положение курсора и изменение положения курсора между кадрами. Эти данные должны быть доступны из пользовательских скриптов, поскольку именно скрипты отвечают за прикладное поведение объектов сцены.

Для трехмерных демонстрационных сцен должно быть предусмотрено управление камерой. Перемещение камеры должно выполняться с помощью клавиш \texttt{W}, \texttt{A}, \texttt{S}, \texttt{D}, \texttt{SPACE} и \texttt{LEFT\_SHIFT}; поворот камеры должен выполняться по смещению мыши. Такой способ управления является привычным для интерактивных трехмерных приложений и позволяет свободно осматривать сцену. Для двумерной интерактивной сцены должно быть предусмотрено получение положения курсора, преобразование координат мыши в координаты мира и обработка нажатий кнопок мыши.

\subsubsection{Требования к системе пользовательских скриптов}

Система должна предоставлять базовый класс сценария поведения. Пользовательский скрипт должен подключаться к сущности как обычный компонент. При запуске сцены система скриптов должна находить такие компоненты и вызывать методы их жизненного цикла.

В пользовательском скрипте разработчик должен иметь возможность получать компоненты текущей сущности, читать пользовательский ввод, изменять трансформацию объекта, прикладывать силы или импульсы к физическому телу, а также обращаться к контексту движка для выполнения действий уровня приложения. Например, в одной из демонстрационных сцен скрипт может менять направление движения кинематической платформы или перезапускать сцену при нажатии клавиши.

\subsubsection{Требования к базовой физической подсистеме}

Физическая подсистема должна обеспечивать базовую обработку физических тел и столкновений. В системе должны поддерживаться динамические, статические и кинематические тела. Динамические тела должны реагировать на гравитацию, скорость, силы и импульсы. Статические тела должны использоваться как неподвижное окружение. Кинематические тела должны перемещаться заданным разработчиком способом и участвовать во взаимодействии с динамическими объектами.

Для проверки столкновений должен использоваться компонент коллайдера. В текущей версии системы основным типом коллайдера является коробочный коллайдер. Он подходит для демонстрации падения кубов, взаимодействия с полом, стенами и движущейся платформой. Требование к физической подсистеме не включает моделирование сложных форм, мягких тел, суставов, транспортных средств и других возможностей полноценных физических движков.

Физическая система должна выполняться после пользовательских скриптов и до рендеринга. Такой порядок позволяет сначала обработать ввод и пользовательскую логику, затем обновить физическое состояние объектов, а после этого отобразить уже измененную сцену.

\subsubsection{Требования к демонстрационным сценам}

В составе программной системы должен быть подготовлен набор демонстрационных сцен, показывающих основные возможности движка в нескольких вариантах использования. Каждая демонстрационная сцена должна запускаться как самостоятельный Java-класс с методом \texttt{main} и использовать публичный программный интерфейс движка для создания приложения, сцены, сущностей, компонентов и систем обработки.

Набор демонстрационных сцен должен включать:
\begin{itemize}
\item \texttt{ShowcaseScene} -- сцену для проверки базового трехмерного рендеринга, работы камеры, направленного освещения, иерархии трансформаций, слоев отображения, отсечения невидимых объектов и отладочных направляющих;
\item \texttt{Interactive2DScene} -- сцену для проверки ортографической камеры, обработки положения мыши, интерактивных объектов, изменения визуального состояния объектов и пользовательских скриптов двумерного интерфейсного типа;
\item \texttt{TexturedLightingScene} -- сцену для проверки материалов с текстурами, процедурного формирования текстур, передачи текстурных данных в OpenGL, направленного освещения и изменения параметров света во время выполнения;
\item \texttt{PhysicsScene} -- сцену для проверки физической подсистемы, свободной камеры, статических, динамических и кинематических тел, коробочных коллайдеров, отладочной визуализации коллайдеров и перезапуска сцены по клавише.
\end{itemize}

Основным демонстрационным примером для проверки физической подсистемы должна выступать сцена \texttt{PhysicsScene}. Она должна содержать статические объекты окружения, динамические физические объекты, падающие под действием гравитации, кинематическую платформу, отладочную сетку и визуализацию коллайдеров. Остальные демонстрационные сцены должны показывать, что архитектура системы не ограничивается физической симуляцией и может использоваться для разных типов интерактивных приложений.

\subsection{Руководство по использованию программной системы}

Данный подраздел описывает порядок работы с разрабатываемой программной системой с точки зрения Java-разработчика. В отличие от конечного пользовательского приложения, движок не предоставляет готовый визуальный интерфейс для создания сцен. Все действия выполняются в исходном коде: разработчик создает приложение, описывает сцену, добавляет объекты и компоненты, после чего запускает программу.

\subsubsection{Подготовка окружения}

Для работы с программной системой необходимо установить JDK версии 25. В качестве внешней библиотеки используется LWJGL версии 3.3.6. В состав подключаемых модулей LWJGL должны входить:
\begin{itemize}
\item \texttt{lwjgl} -- базовый модуль LWJGL;
\item \texttt{lwjgl-glfw} -- модуль для создания окна, OpenGL-контекста и обработки ввода;
\item \texttt{lwjgl-opengl} -- модуль для доступа к функциям OpenGL;
\item \texttt{lwjgl-stb} -- модуль для загрузки изображений и текстур.
\end{itemize}

Проект не использует Maven или Gradle. Поэтому разработчик может работать с ним двумя способами. Первый способ заключается в открытии проекта в IntelliJ IDEA и ручном добавлении JAR-файлов LWJGL в зависимости проекта. Второй способ заключается в ручной компиляции исходных файлов с помощью \texttt{javac} и запуске нужного класса через \texttt{java}.

При настройке проекта через IntelliJ IDEA необходимо выполнить следующие действия:
\begin{enumerate}
\item Открыть каталог проекта как обычный Java-проект.
\item Указать установленный JDK 25 в качестве SDK проекта.
\item Добавить JAR-файлы LWJGL 3.3.6 и JAR-файлы нативных библиотек для целевой операционной системы в раздел зависимостей проекта.
\item Убедиться, что рабочим каталогом запуска является корневой каталог проекта, в котором расположен каталог \texttt{res}.
\item Создать конфигурации запуска для одного или нескольких демонстрационных классов: \texttt{common.showcase.ShowcaseScene}, \texttt{common.showcase.Interactive2DScene}, \texttt{common.showcase.TexturedLightingScene}, \texttt{common.showcase.PhysicsScene}.
\end{enumerate}

Для ручной компиляции можно использовать следующий общий порядок команд. Команды предполагают, что все JAR-файлы LWJGL помещены в каталог \texttt{lib}, а компиляция выполняется из корня проекта. В примерах ниже запускается сцена \texttt{PhysicsScene}; для запуска других демонстрационных сцен достаточно заменить имя запускаемого класса на \texttt{ShowcaseScene}, \texttt{Interactive2DScene} или \texttt{TexturedLightingScene} из пакета \texttt{common.showcase}.

\begin{lstlisting}[caption={Пример ручной компиляции проекта в Windows}]
dir /s /b src\*.java test\*.java > sources.txt
javac -encoding UTF-8 -d out -cp "lib/*" @sources.txt
java -cp "out;lib/*" common.showcase.PhysicsScene
\end{lstlisting}

Для Linux и macOS отличается только разделитель элементов classpath. Вместо точки с запятой используется двоеточие.

\begin{lstlisting}[caption={Пример ручной компиляции проекта в Linux и macOS}]
find src test -name "*.java" > sources.txt
javac -encoding UTF-8 -d out -cp "lib/*" @sources.txt
java -cp "out:lib/*" common.showcase.PhysicsScene
\end{lstlisting}

Если нативные библиотеки LWJGL не находятся автоматически через JAR-файлы, их необходимо указать через параметр \texttt{java.library.path}. Однако при использовании набора библиотек, полученного через конфигуратор LWJGL вместе с native JAR для соответствующей операционной системы, обычно достаточно добавить все JAR-файлы в classpath.

\subsubsection{Создание приложения}

Точкой входа в интерактивное приложение является обычный метод \texttt{main}. В нем создается конфигурация движка, экземпляр приложения, загружается стартовая сцена и запускается основной цикл. Пример минимальной структуры запуска приведен в листинге 2.3.

\begin{lstlisting}[language=Java,caption={Создание приложения и запуск сцены}]
public static void main(String[] args) {
    EngineConfig config = new EngineConfig()
        .clearColor(0.05f, 0.07f, 0.08f, 1f)
        .autoUpdateCameraAspect(true)
        .debugOpenGLErrors(false);

    Application app = new Application("interactive scene", config);
    app.getWindow().setCursorDisabled();
    app.context().sceneManager().load(createScene());
    app.run();
}
\end{lstlisting}

В данном примере объект \texttt{EngineConfig} определяет параметры запуска. Метод \texttt{clearColor} задает цвет очистки окна перед отрисовкой кадра. Параметр \texttt{autoUpdateCameraAspect} включает автоматическое обновление соотношения сторон камеры при изменении размеров окна. Параметр \texttt{debugOpenGLErrors} позволяет отключить или включить дополнительные проверки ошибок OpenGL.

После создания приложения в трехмерных демонстрационных сценах может отключаться отображение курсора, так как они используют свободное управление камерой мышью. В двумерной интерактивной сцене курсор, наоборот, остается видимым для обработки кликов по объектам. Затем через менеджер сцен загружается сцена, созданная методом \texttt{createScene}, и вызывается метод \texttt{run}. После этого управление переходит основному циклу приложения.

\subsubsection{Создание сцены и добавление систем}

Сцена создается как экземпляр класса \texttt{Scene}. До добавления сущностей в нее необходимо подключить системы обработки. Порядок подключения систем важен, потому что каждая система выполняется в порядке добавления. В листинге ниже приведен пример настройки систем для физической демонстрационной сцены.

\begin{lstlisting}[language=Java,caption={Создание сцены и подключение систем}]
private static Scene createScene() {
    Scene scene = new Scene();

    scene.addSystem(new ScriptSystem());

    scene.addSystem(new PhysicsSystem()
        .gravity(0f, -9.81f, 0f)
        .fixedTimeStep(1f / 60f)
        .maxSubSteps(5)
        .velocityIterations(12)
        .positionIterations(10)
        .debugDrawColliders(true));

    scene.addSystem(new RenderSystem()
        .ambientColor(new Vec4(0.18f, 0.20f, 0.23f, 1f)));

    scene.addSystem(new DebugRenderSystem());

    return scene;
}
\end{lstlisting}

Сначала добавляется \texttt{ScriptSystem}, поскольку пользовательские скрипты могут изменять состояние объектов, читать ввод, прикладывать силы или запрашивать смену сцены. После этого добавляется \texttt{PhysicsSystem}, которая обновляет физические тела и коллайдеры. Затем выполняется \texttt{RenderSystem}, получающая уже обновленные трансформации объектов. Последней добавляется \texttt{DebugRenderSystem}, отвечающая за вывод отладочной геометрии.

\subsubsection{Добавление источника света}

Для базового освещения сцены используется компонент \texttt{DirectionalLight}. Он задает направление света, цвет и интенсивность. Такой источник света не имеет положения в пространстве и имитирует удаленный источник, например солнце.

\begin{lstlisting}[language=Java,caption={Создание направленного источника света}]
scene.addEntity(new Entity("sun")
    .addComponent(new DirectionalLight()
        .direction(-0.4f, -1f, -0.25f)
        .color(new Vec4(1f, 0.94f, 0.82f, 1f))
        .intensity(1.2f)));
\end{lstlisting}

Свет оформляется как обычная сущность с компонентом. Такой подход соответствует общей компонентной модели: даже служебные объекты сцены могут быть представлены через сущности и компоненты.

\subsubsection{Создание камеры}

Для отображения сцены необходима сущность с компонентом \texttt{Camera}. Чтобы камера имела положение и поворот, к сущности добавляется компонент \texttt{Transform}. Для управления камерой используется скрипт \texttt{FreecamController}.

\begin{lstlisting}[language=Java,caption={Создание свободной камеры}]
scene.addEntity(new Entity("camera")
    .addComponent(new Transform()
        .position(0f, 4f, 9f)
        .rotation(Quaternion.euler(-22f, 0f, 0f)))
    .addComponent(new Camera()
        .fov(60f)
        .near(0.1f)
        .far(80f)
        .cullingMask(LAYER_WORLD | LAYER_PHYSICS)
        .priority(10))
    .addComponent(new FreecamController()
        .speed(7f)
        .sensitivity(18f)));
\end{lstlisting}

Компонент \texttt{Camera} задает угол обзора, ближнюю и дальнюю плоскости отсечения, маску видимости и приоритет. Приоритет нужен в случае, если в сцене присутствует несколько камер. Скрипт \texttt{FreecamController} считывает ввод пользователя и изменяет трансформацию камеры.

\subsubsection{Создание статического объекта}

Статический объект окружения можно создать без компонента \texttt{Rigidbody}. Для пола физической демонстрационной сцены достаточно задать трансформацию, коробочный коллайдер и компонент отображения меша.

\begin{lstlisting}[language=Java,caption={Создание пола сцены}]
scene.addEntity(new Entity("floor")
    .addComponent(new Transform()
        .position(0f, -0.1f, 0f)
        .scale(8f, 0.2f, 8f))
    .addComponent(new BoxCollider()
        .halfExtents(0.5f, 0.5f, 0.5f)
        .friction(1f)
        .restitution(0f))
    .addComponent(new MeshRenderer()
        .mesh(Meshes.cube())
        .material(new Material()
            .color(new Vec4(0.18f, 0.24f, 0.22f, 1f)))
        .layer(LAYER_WORLD)
        .boundsRadius(6f)));
\end{lstlisting}

В данном примере куб используется как масштабируемый параллелепипед. Компонент \texttt{Transform} растягивает его до размеров пола. Компонент \texttt{BoxCollider} задает область столкновения. Компонент \texttt{MeshRenderer} указывает геометрию, материал, слой отображения и радиус границ, используемый при обработке видимости.

\subsubsection{Создание динамического физического тела}

Динамический объект отличается от статического тем, что к нему добавляется компонент \texttt{Rigidbody}. Этот компонент хранит тип тела, массу, скорость, демпфирование, настройки гравитации и другие физические параметры.

\begin{lstlisting}[language=Java,caption={Создание динамического куба}]
scene.addEntity(new Entity("falling-box")
    .addComponent(new Transform()
        .position(0f, 5f, 0f))
    .addComponent(new Rigidbody()
        .mass(1f)
        .velocity(1f, 0f, 0f)
        .linearDamping(0.05f)
        .angularDamping(0.08f))
    .addComponent(new BoxCollider()
        .halfExtents(0.5f, 0.5f, 0.5f)
        .friction(0.9f)
        .restitution(0.1f))
    .addComponent(new MeshRenderer()
        .mesh(Meshes.cube())
        .material(new Material()
            .color(new Vec4(0.25f, 0.68f, 1f, 1f)))
        .layer(LAYER_PHYSICS)
        .boundsRadius(1f)));
\end{lstlisting}

После запуска сцены физическая система будет применять к такому объекту гравитацию, учитывать скорость, выполнять проверку столкновений с коллайдерами и изменять компонент \texttt{Transform}. Рендеринг будет использовать уже обновленное положение объекта.

\subsubsection{Создание кинематического объекта}

Кинематическое тело не падает под действием гравитации, но может перемещаться с заданной скоростью и взаимодействовать с динамическими телами. Такой объект удобно использовать для движущихся платформ и элементов окружения.

\begin{lstlisting}[language=Java,caption={Создание кинематической платформы}]
scene.addEntity(new Entity("kinematic-platform")
    .addComponent(new Transform()
        .position(0f, 0.35f, 3.1f)
        .scale(1.8f, 0.25f, 0.7f))
    .addComponent(new Rigidbody()
        .type(PhysicsBodyType.KINEMATIC)
        .velocity(0.9f, 0f, 0f)
        .useGravity(false))
    .addComponent(new BoxCollider()
        .halfExtents(0.5f, 0.5f, 0.5f)
        .friction(0.9f))
    .addComponent(new MeshRenderer()
        .mesh(Meshes.cube())
        .material(new Material()
            .color(new Vec4(1f, 0.9f, 0.28f, 1f)))
        .layer(LAYER_PHYSICS)
        .boundsRadius(2f))
    .addComponent(new PlatformBounceScript(-2.4f, 2.4f)));
\end{lstlisting}

В конце к сущности добавляется пользовательский скрипт \texttt{PlatformBounceScript}. Он следит за положением платформы и меняет направление скорости при достижении границ движения.

\subsubsection{Создание пользовательского скрипта}

Пользовательский скрипт создается наследованием от базового класса \texttt{Script}. В методе \texttt{init} обычно выполняется получение ссылок на компоненты сущности, а в методе \texttt{loop} описывается поведение, выполняемое каждый кадр.

\begin{lstlisting}[language=Java,caption={Пример пользовательского скрипта движения платформы}]
private static final class PlatformBounceScript extends Script {
    private final float minX;
    private final float maxX;
    private Rigidbody body;
    private Transform transform;

    private PlatformBounceScript(float minX, float maxX) {
        this.minX = minX;
        this.maxX = maxX;
    }

    @Override
    public void init() {
        body = entity().getComponent(Rigidbody.class);
        transform = entity().getComponent(Transform.class);
    }

    @Override
    public void loop(float dtime) {
        if (body == null || transform == null) return;

        Vec3 position = transform.position();
        Vec3 velocity = body.velocity();

        if ((position.x <= minX && velocity.x < 0f) ||
            (position.x >= maxX && velocity.x > 0f)) {
            body.velocity(-velocity.x, velocity.y, velocity.z);
        }
    }
}
\end{lstlisting}

Данный пример показывает типовой способ расширения поведения объектов. Скрипт не изменяет устройство движка, не требует создания новой системы и работает как компонент конкретной сущности. Благодаря этому прикладная логика остается отделенной от ядра программной системы.

\subsubsection{Использование пользовательского ввода}

Для чтения пользовательского ввода из скриптов используется класс \texttt{Input}. Например, в физической демонстрационной сцене клавиша \texttt{R} может использоваться для перезапуска сцены.

\begin{lstlisting}[language=Java,caption={Пример обработки нажатия клавиши}]
private static final class PhysicsSceneControlScript extends Script {
    @Override
    public void loop(float dtime) {
        if (!Input.getKeyDown(KeyCode.R)) return;

        EngineContext context = EngineContext.current();
        context.window().setCursorDisabled();
        context.sceneManager().requestReplace(createScene());
    }
}
\end{lstlisting}

В этом примере используется однократное нажатие клавиши. При нажатии \texttt{R} получается текущий контекст движка, скрывается курсор и запрашивается замена активной сцены. Замена выполняется не немедленно, а через менеджер сцен, что предотвращает изменение структуры приложения в небезопасный момент цикла обновления.

\subsubsection{Создание отладочной геометрии}

Для наглядной проверки работы сцены может использоваться отладочный рендеринг. Например, можно каждый кадр выводить сетку на плоскости, чтобы лучше видеть движение физических объектов.

\begin{lstlisting}[language=Java,caption={Пример вывода отладочной сетки}]
private static final class PhysicsGridScript extends Script {
    @Override
    public void loop(float dtime) {
        for (int i = -4; i <= 4; i++) {
            DebugRenderer.line(
                new Vec3(i, 0f, -4f),
                new Vec3(i, 0f, 4f),
                new Vec4(0.18f, 0.25f, 0.28f, 1f));

            DebugRenderer.line(
                new Vec3(-4f, 0f, i),
                new Vec3(4f, 0f, i),
                new Vec4(0.18f, 0.25f, 0.28f, 1f));
        }
    }
}
\end{lstlisting}

Отладочная геометрия не является частью постоянной модели сцены. Она создается на кадр и используется для проверки положения объектов, границ коллайдеров, ориентации сцены и поведения физической подсистемы.

\subsubsection{Общий порядок создания сцены}

С учетом приведенных примеров общий порядок разработки интерактивной сцены выглядит следующим образом:
\begin{enumerate}
\item Создается конфигурация движка и экземпляр приложения.
\item Создается объект сцены.
\item В сцену добавляются системы обработки в требуемом порядке.
\item При необходимости в сцену добавляется источник света.
\item В сцену добавляется камера: перспективная для трехмерной сцены или ортографическая для двумерной сцены.
\item Создаются сущности окружения, интерактивные объекты или физические тела в зависимости от назначения сцены.
\item При необходимости добавляются пользовательские скрипты поведения.
\item Сцена загружается через менеджер сцен.
\item Запускается основной цикл приложения.
\end{enumerate}

Такой порядок хорошо подходит для учебных и демонстрационных проектов, потому что вся сцена описывается в одном или нескольких Java-классах. Разработчик сразу видит, какие системы включены в сцену, какие сущности созданы, какие компоненты назначены объектам и какой код отвечает за поведение.

\subsection{Требования к программному интерфейсу}

Программный интерфейс системы должен быть доступен из Java-кода и не должен требовать от разработчика прямой работы с низкоуровневыми вызовами OpenGL для типовых операций создания сцены. Разработчик должен использовать готовые классы приложения, сцены, сущностей, компонентов, систем, материалов, мешей, текстур, ввода и времени.

Публичные классы должны поддерживать цепочечную настройку объектов, когда метод изменения параметра возвращает сам объект. Такой стиль используется при создании конфигурации, компонентов, материалов и систем. Он уменьшает объем вспомогательного кода и делает описание сцены более компактным.

Программный интерфейс должен позволять:
\begin{itemize}
\item создавать приложение с заданным заголовком окна и конфигурацией;
\item загружать и заменять активную сцену;
\item добавлять в сцену системы обработки;
\item создавать сущности и назначать им компоненты;
\item получать компоненты сущности по типу;
\item создавать пользовательские скрипты поведения;
\item получать состояние клавиатуры и мыши;
\item загружать и использовать текстуры, материалы и шейдеры;
\item создавать физические тела и коробочные коллайдеры;
\item выводить отладочную геометрию при разработке сцены.
\end{itemize}

При этом интерфейс должен оставаться достаточно простым для учебного использования. Разработчик должен иметь возможность создать первую интерактивную сцену без предварительного изучения сложного визуального редактора, системы сборки, специализированного языка описания сцен или внешнего формата проекта.

\subsection{Нефункциональные требования}

\subsubsection{Требования к программному обеспечению}

Для разработки и запуска программной системы требуется JDK 25. В качестве внешней библиотеки используется LWJGL 3.3.6. Обязательными модулями LWJGL являются GLFW, OpenGL и STB. Система должна запускаться на настольных операционных системах Windows, Linux и macOS при наличии соответствующих нативных библиотек LWJGL и графического драйвера.

Так как система использует OpenGL и шейдеры GLSL версии 330 core, целевая рабочая станция должна поддерживать OpenGL 3.3 Core Profile или более новую совместимую версию. При отсутствии подходящего графического драйвера корректная работа графической подсистемы не гарантируется.

\subsubsection{Требования к аппаратному обеспечению}

Аппаратные требования программной системы определяются требованиями JDK 25, LWJGL, используемых нативных библиотек и графического драйвера. Для работы демонстрационных сцен необходим персональный компьютер или ноутбук с поддержкой OpenGL 3.3 Core Profile, клавиатурой и мышью. Отдельные численные требования к объему оперативной памяти, производительности процессора и видеосистемы в рамках технического задания не задаются, поскольку демонстрационные сцены не ориентированы на предельную нагрузку и используются для проверки архитектурных возможностей движка.

\subsubsection{Требования к надежности}

Программная система должна предсказуемо обрабатывать некорректное использование основных объектов. При передаче недопустимых значений в конфигурацию, компоненты или системы должны возникать диагностируемые исключения. Например, не допускается создание сущности с пустым идентификатором, добавление пустого компонента, установка отрицательного фиксированного шага времени или установка неположительной массы физического тела.

Система должна корректно освобождать ресурсы при завершении приложения. После выхода из основного цикла должны быть уничтожены активные сцены, компоненты, системы, окно, графические ресурсы и контекст выполнения. Это требование особенно важно для приложений, использующих нативные ресурсы OpenGL и LWJGL.

Изменение состава сцены во время обновления должно выполняться безопасно. Если сущность, компонент или система добавляются или удаляются из пользовательского скрипта, операция должна быть отложена до момента, когда изменение структуры сцены не нарушит обход систем и сущностей.

\subsubsection{Требования к сопровождаемости}

Исходный код программной системы должен иметь модульную структуру. Классы ядра компонентной модели, компоненты, графическая подсистема, скрипты, физическая подсистема и демонстрационные сцены должны быть разделены по пакетам. Такое разделение необходимо для того, чтобы разработчик мог изучать и изменять отдельные части движка без просмотра всего проекта сразу.

Демонстрационные сцены должны использовать публичный программный интерфейс движка и выступать как практические примеры его применения. Это позволяет использовать их не только для проверки работоспособности, но и как основу пользовательской документации.

\subsubsection{Требования к безопасности}

Разрабатываемая программная система является локальной программной основой и не предусматривает сетевого взаимодействия, хранения учетных записей, обработки персональных данных и разграничения пользовательских ролей. Поэтому требования информационной безопасности ограничиваются корректной обработкой внешних файлов ресурсов и предотвращением аварийного завершения при ошибках загрузки.

При работе с ресурсами должны быть диагностируемыми ошибки чтения файлов, компиляции шейдеров, загрузки текстур и создания графических объектов. Разработчик должен иметь возможность понять причину ошибки по сообщению исключения или отладочному выводу.

\subsection{Требования к оформлению документации}

Программная документация должна включать материалы, необходимые для понимания назначения, архитектуры, порядка сборки, запуска и использования программной системы. В составе пояснительной записки должны быть представлены анализ предметной области, техническое задание, технический проект и рабочий проект.

В техническом задании должны быть описаны назначение разработки, задачи, требования к программной системе, функциональные сценарии и порядок использования движка разработчиком. В техническом проекте должны быть раскрыты выбранные средства разработки, архитектура, структура пакетов, основные подсистемы и проектные решения. В рабочем проекте должны быть приведены спецификации основных классов, описание тестирования и результаты проверки работоспособности программной системы.

Текстовые материалы, рисунки, таблицы, листинги и графический материал должны оформляться в соответствии с требованиями, применяемыми к выпускным квалификационным работам по направлению подготовки 09.03.04 «Программная инженерия».
